%-----------------------------------------------------------------------
% MIXCOMPILE: A Program-Feature-Driven Adaptive Method for Diversified
%             Compilation Pass Selection
% Version: v5-slim (2026-10-10), condensed from v5 (2026-10-09)
%   Streamlined revision:
%     - duplicated disclosure narratives (E6 / T1 / dc incident) stated
%       once, in the section that owns the evidence
%     - "Design Challenges" section merged into the method overview;
%       the pipeline is now described once instead of three times
%     - "Summary of Findings" section folded into the Conclusion
%     - E3/E5/E8/T5 compressed; related-work LLM subsection no longer
%       duplicates the introduction's attack-side survey
%     - build-time and decompilation-time figures dropped (captions
%       restated the tables verbatim)
%     - templated phrasing and repetitive honesty disclaimers reduced
% For Springer Nature Journal (sn-jnl article template)
%-----------------------------------------------------------------------

\documentclass[pdflatex,sn-mathphys-num]{sn-jnl}

%%%% Standard Packages
\usepackage{graphicx}%
\usepackage{multirow}%
\usepackage{amsmath,amssymb,amsfonts}%
\usepackage{amsthm}%
\usepackage{mathrsfs}%
\usepackage[title]{appendix}%
\usepackage{xcolor}%
\usepackage{textcomp}%
\usepackage{booktabs}%
\usepackage{algorithm}%
\usepackage{algorithmicx}%
\usepackage{algpseudocode}%
\usepackage{listings}%

% TODO marker for pending experimental numbers (visible placeholders)
\newcommand{\todonote}[1]{\textcolor{red}{[TODO: #1]}}

%%%% Theorem Styles
\theoremstyle{thmstyleone}%
\newtheorem{theorem}{Theorem}%
\newtheorem{proposition}[theorem]{Proposition}%

\theoremstyle{thmstyletwo}%
\newtheorem{example}{Example}%
\newtheorem{remark}{Remark}%

\theoremstyle{thmstylethree}%
\newtheorem{definition}{Definition}%

\raggedbottom

\begin{document}

%%% Title
\title[Program-Feature-Driven Adaptive Diversified Pass Selection]{Program-Feature-Driven Adaptive Diversified Pass Selection}

%%% Authors
\author[1]{\fnm{Dingyu} \sur{Wang}}

\author*[2]{\fnm{Junchao} \sur{Wang}}\email{wangjunc@ustc.edu.cn}
\equalcont{These authors contributed equally to this work.}

\author*[3]{\fnm{Fan} \sur{Zhang}}\email{ffzhang@fudan.edu.cn}

\author[4]{\fnm{Guangwei} \sur{Xie}}

\author[3]{\fnm{Mingqian} \sur{Sun}}

%%% Affiliations
\affil[1]{\orgdiv{School of Cyberspace Security}, \orgname{Zhengzhou University}, \orgaddress{\city{Zhengzhou}, \postcode{450002}, \country{China}}}

\affil*[2]{\orgdiv{School of Cyberspace Security}, \orgname{Information Engineering University}, \orgaddress{\city{Zhengzhou}, \postcode{450001}, \country{China}}}

\affil*[3]{\orgdiv{Institute of Big Data}, \orgname{Fudan University}, \orgaddress{\city{Shanghai}, \postcode{200433}, \country{China}}}

\affil[4]{\orgname{Shanghai Hongzhen Information Technology Co., Ltd.}, \orgaddress{\city{Shanghai}, \postcode{201112}, \country{China}}}

%%% Abstract
\abstract{
Compiler-level software diversification denies adversaries the amortization of a
one-time reverse-engineering investment across identical deployments, but the
economics of this defense is deteriorating: large language models (LLMs) collapse
the marginal cost of binary adaptation, while existing obfuscating compilers apply
uniform pipelines, offer no quantitative control over the security--performance
trade-off, and expose no interface that automated compilation agents could invoke.
We present MIXCOMPILE, an LLVM~17-based framework that makes each deployed binary
an independent, cost-bounded sampling decision. Its foundation is a
five-dimensional cost model (security gain, diversity gain, runtime, size,
correctness risk) that gives LLM-based compilation agents a structured,
quantitative basis for reasoning about security--performance trade-offs. Its
architecture decouples selection from execution through LLVM metadata annotations
and exposes a machine-readable interface---strategy space, per-decision audit log,
and stable invocation surface---through which an LLM can plan, invoke, and verify
compilation strategies; on this interface, LLM planning preserves semantics in
98.9\% of cases, versus 63.2\% for direct LLM source rewriting. Its selection
mechanism operates at function granularity: single-traversal feature extraction
shapes each program's strategy by its own structure and makes recompilations
structurally distinct, targeting the amortization economics of LLM-assisted
adversaries. On a frozen 44/43/43 program split, the pre-registered weight
configuration dominates 28 of 30 random-weight candidates, 4 of which silently
corrupt correctness; all 30 seeds of all 12 security-subset programs pass semantic
verification, and the resulting 30-variant pools are 100\% unique at the
\texttt{.text} level. On real workloads, MIXCOMPILE stays within 1--13\% of LLVM
throughput on llama.cpp, completes the full build in 2.0$\times$ the LLVM build
time (uniform full obfuscation exceeds a three-hour limit without producing a
binary), and increases Ghidra decompilation cost to 8.11$\times$ LLVM, with
function-recovery F1 reduced from 0.903 to 0.708 and classical
function-localization accuracy reduced to 0.163 (unobfuscated $\approx$1.0).
}
%%% Keywords
\keywords{software diversity, obfuscation pass selection, five-dimensional cost model, program feature extraction, LLM-based compilation, LLVM compiler}


\maketitle

%%%%%%%%%%%%%%%%%%%%%%%%%%%%%%%%%%%%%%%%%%%%%%%%%%%%%%%
%%% Main text
%%%%%%%%%%%%%%%%%%%%%%%%%%%%%%%%%%%%%%%%%%%%%%%%%%%%%%%

%=======================================================
\section{Introduction}
%=======================================================

Modern software is typically deployed across cyberspace as millions of
structurally identical binaries. For an adversary this is an economic
advantage: the cost of understanding one binary is paid once and amortized over
every deployment of the same build. Compiler-level software diversification
has been studied for over a decade to break this
amortization~\cite{bib2,franz2010unibus,homescu2013librando,crane2015readactor}:
if every deployment is an independent structural draw, the attacker must
re-adapt per target, and defense value equals the marginal increase in
adaptation cost. The same argument underlies moving-target and
endogenous-security architectures in system design~\cite{bib1}.

The attack side of this economics has now changed qualitatively. Automated
adaptation is not new: just-in-time code-reuse attacks showed in 2013 that a
single memory disclosure suffices to re-assemble an attack against randomized
code at runtime~\cite{snow2013jitrop}. What is new is that large language
models (LLMs) perform the same work at near-zero
cost~\cite{pearce2022asleep,perry2023users}: they generate functional variants
that fully evade four state-of-the-art detection tools in the hardware IP
domain~\cite{gohil2025llmpirate}; autonomous attack agents complete individual
attacks for under \$15 with 37/40 attempted tasks
succeeding~\cite{singer2026incalmo}; and obfuscation-aware fine-tuning
recovers up to 16.9 percentage points of instruction-following accuracy on
commercial virtualization-based protectors~\cite{wang2026oasif}. A defense
whose value is measured in attacker adaptation cost must therefore be
re-examined when adaptation becomes cheap.

Against this background, existing obfuscating compilers leave three problems
open. First, uniform pipelines of the Obfuscator-LLVM
family~\cite{bib8,bib9,bib10,bib11,bib12} apply the same pass set to every
function of every program: there is no selection mechanism and no cost
accounting, and repeated compilations of the same source produce identical
binaries, which re-enables the very amortization that diversification is meant
to break. Second, randomized compilation
frameworks~\cite{homescu2013librando,crane2015readactor,koo2018ccr,franzen2023randcompile}
introduce seed-driven variation, but their randomization is cost-blind---no
component estimates what a transformation will cost before applying it---and
their threat model targets memory-disclosure and code-reuse attacks rather than
amortized reverse engineering. Third, learning-based selective
obfuscation~\cite{bib4,wei2025reinobf} optimizes a single objective, typically
evasion of a known similarity detector, with at most two cost dimensions
folded into the reward and no independent semantic verification of the output.
The newest techniques do not close the gap: LLMs that obfuscate source code
directly achieve at best an 81\% transformation pass rate and tend to
\emph{simplify} the code rather than obscure
it~\cite{detomasi2025simplicity}, and LLM-guided compilation
workflows~\cite{li2026autopass} operate exclusively on the performance axis
without an interface for planning or invoking security-directed strategies.

To address these problems, we present MIXCOMPILE, an LLVM~17-based framework
built on one insight: \emph{make every deployed binary an independent sampling
decision under an explicit cost budget}, and make that decision computable,
auditable, and invocable. The framework rests on three coupled components. A
five-dimensional cost model (security gain, diversity gain, runtime cost, size
cost, correctness risk) scores each of seven obfuscation passes against the
feature distribution of the program, turning the security--performance
trade-off into structured, quantitative data that LLM-based compilation agents
can reason over instead of prompt-level intuition. A decoupled
selection--execution architecture emits every decision as an LLVM metadata
annotation and exposes a machine-readable interface through which external
callers, including LLM-based compilation
workflows, can plan, invoke, and verify compilation strategies while the
constructive safeguards remain inside
the compiler. A function-granularity selection mechanism extracts control-flow
and instruction-composition features in a single traversal, so that each
program's strategy is shaped by its own structure and boundary randomization
makes recompilations structurally distinct, directly targeting the amortization
economics of LLM-assisted adversaries.

Formally, the paper asks a narrow, empirically decidable question at the
compiler level: \emph{can a compiler automatically select, for each program, a
combination of obfuscation passes that (i) measurably increases the cost of
automated binary analysis, (ii) keeps runtime, code-size, and correctness risk
under explicit quantitative control, and (iii) produces mutually distinct
variants across repeated compilations of the same source, so that an
adversary's adaptation cannot be amortized across deployments?} The threat
model is explicit: the adversary is an automated binary-analysis toolchain or
an LLM-assisted analyst operating on deployed binaries. The defense goal is
to raise the per-target adaptation cost, and we do not claim unbreakability
against human experts, cryptographic security, or immunity to a sufficiently
resourced attacker.

The contributions of this paper are threefold.
\begin{itemize}
\item \textbf{A five-dimensional, LLM-invocable cost model for obfuscation pass
  selection}, giving LLM-based compilation agents a structured, quantitative basis
  for reasoning about security--performance trade-offs instead of prompt-level
  intuition.
\item \textbf{A decoupled, modular selection--execution architecture.} Selection
  decisions are emitted as LLVM metadata and exposed through a machine-readable
  interface---strategy space, per-decision audit log, and stable invocation
  surface---through which an LLM can plan, invoke, and verify compilation
  strategies; new passes are admitted by declaration alone.
\item \textbf{A function-granularity adaptive selection mechanism for the LLM
  era.} Single-traversal feature extraction shapes each program's strategy by its
  own structure and makes recompilations structurally distinct, targeting the
  amortization economics of LLM-assisted adversaries.
\end{itemize}

The remainder of this paper is organized as follows. Section~2 introduces
related work and locates the vacancy in a coordinate system of existing
mechanisms. Section~3 presents the design, mapping four design challenges
onto the three components. Section~4 reports the experimental evaluation.
Section~5 concludes.

%=======================================================
\section{Related Work}
%=======================================================

\subsection{Software Diversity and Attack Amortization}

The security value of software diversity rests on a single economic observation: an
attack engineered against one binary should not transfer for free to another. Unibus
used diversity through redundancy to survive arbitrary software
bugs~\cite{franz2010unibus}; the SoK of Larsen et al.~\cite{bib2} systematized
automated software diversity and identified \emph{amortization} as the attacker's core
economic advantage. Compiler-level instantiations include Librando's transparent
randomization of JIT-compiled applications~\cite{homescu2013librando}, Readactor's code
randomization resilient to memory disclosure~\cite{crane2015readactor},
compiler-assisted code-reuse disruption~\cite{koo2018ccr}, and randomized per-build
removal of forensic gadgets from the Linux kernel~\cite{franzen2023randcompile}.
Reproducible-builds research notes the opposite trade-off, identical binaries aiding
auditing and supply-chain verification~\cite{fourn2023reproducible}, which sharpens the
point that diversity is a \emph{deployment} decision rather than a universal default:
a diversity mechanism earns its cost only when it measurably raises adaptation cost.
These frameworks randomize \emph{structure} but do not decide \emph{what to
transform}; no component estimates, before applying a transformation, what it will
cost in runtime, size, or correctness risk.

\subsection{Obfuscation Compilation and Selective Obfuscation}

Obfuscation theory and taxonomy date to Collberg et al.~\cite{bib6};
Schrittwieser et al.~\cite{bib7} measured how long obfuscation actually holds against
modern analysis, and automated deobfuscation is already practical on native
code~\cite{bib5}. In engineering practice, Obfuscator-LLVM (OLLVM)~\cite{bib8} fixed
the three-pass canon (instruction substitution, bogus control flow, control-flow
flattening) applied uniformly, and successors strengthen individual transforms:
extended instruction replacement~\cite{bib9}, nested-switch and in-degree control-flow
obfuscation~\cite{bib10}, code rearrangement~\cite{bib11}, and fragment
fusion~\cite{bib12}. Optimization-space obfuscation~\cite{bib13} and genetic-algorithm
search over LLVM passes~\cite{bib14} show that the compiler's own machinery can
obstruct analysis, and evaluation-methodology work~\cite{bib15} warns that
obfuscation papers often skip resistance measurement. What unifies this line is the
absence of a selection mechanism: every function receives the same transformation,
however improved.

Selective obfuscation is emerging. Wang et al.~\cite{bib4} learn obfuscation sequences
with reinforcement learning against a similarity detector; ReinObf~\cite{wei2025reinobf}
optimizes lightweight block-level obfuscation against a known ensemble of binary-code
similarity detectors, with two cost terms folded into its reward. Both optimize a
single external objective, verify semantics post hoc at most, and produce one sequence
rather than a variant pool. Industrial compile-time isomerization~\cite{he2024isomerization}
generates heterogeneous executables for redundant deployment but publishes no
selection algorithm or cost accounting. At the transform level,
Polaris~\cite{zhou2026polaris} pushes obfuscation below IR to the machine IR to resist
decompiler mid-end analyses, at $3.39$--$5.95\times$ size overhead and without
per-program adaptation, which its authors list as future work. Transform research
keeps deepening, and selection and cost control remain open. These two concerns are
orthogonal to transform quality: deeper transforms enlarge the candidate pool that a
selector must reason about.

\subsection{LLM-Era Defenses}

The arrival of LLMs has begun to reshape the defense side as well, though
imprecisely. LLM-based obfuscation achieves at best an 81\% transformation
pass rate on real functions, and the transformed code is consistently
\emph{simplified} rather than obscured~\cite{detomasi2025simplicity}.
LLM-generated program diversity improves reliability in N-version settings
but operates at source level and small scale~\cite{ron2026galapagos}.
LLM-guided compilation agents~\cite{li2026autopass} demonstrate that an LLM
can drive pass ordering when the compiler is opened up and provides grounded
feedback, but only for performance. The compiler-level response to cheap
adversarial adaptation is missing. This paper fills it with diversified,
cost-bounded, semantically safe compilation behind a machine-readable
strategy interface that LLM workflows can plan over, invoke, and verify.

\subsection{Compiler Pass Selection and Cost Modeling}

Compiler testing established that pass combinations form an explorable
space~\cite{bib16,bib17,bib18,bib19}, and machine-learning pass selection for
\emph{performance} is mature: reinforcement-learned phase ordering~\cite{bib20},
coreset-based sequence prediction~\cite{bib21}, and compiler world
models~\cite{bib22}. Program features predict optimization outcomes~\cite{bib25,bib26},
and binary-code similarity provides the measurement
instruments~\cite{bib23,bib24,bib27}. MIXCOMPILE transfers this methodology from the
performance axis to the security axis, where the objective space is different in kind:
security and diversity gains are estimates, runtime and size are measurable costs, and
correctness risk is a constraint that must never be traded away silently. The model
accordingly treats correctness as a first-class dimension rather than a post-hoc
test.

\subsection{Position: A Coordinate System and Its Vacancy}
\label{sec:coordsys}

Table~\ref{tab:coordsys} arranges the mechanisms most relevant to this paper along six
axes: transformation layer, selection mechanism, explicit cost model, semantic
guarantee, diversity product, and threat model. The vacancy is structural, not
rhetorical: no verified framework spans all six axes. Polaris lists adaptive tuning
as future work~\cite{zhou2026polaris}; ReinObf optimizes a single detector
objective~\cite{wei2025reinobf}; randomization frameworks are
cost-blind~\cite{bib2,franzen2023randcompile}; OLLVM applies uniformly~\cite{bib8}.
MIXCOMPILE occupies the vacancy: it does not propose new transforms (transform quality
is complementary and orthogonal) but supplies the missing selection, cost, safety, and
diversity layers on top of a standard pass pool.

\begin{table*}[!t]
\centering
\footnotesize
\caption{Coordinate system of related mechanisms. ``---''~means the axis is absent in
the published mechanism. MIXCOMPILE row in bold. Polaris overhead figures are from its
synthetic (Csmith) evaluation~\cite{zhou2026polaris}.}
\label{tab:coordsys}
\begin{tabular}{lllllll}
\toprule
Approach & Layer & Selection mechanism & Cost model & Semantic guarantee & Diversity product & Threat model \\
\midrule
OLLVM~\cite{bib8} & IR & none (uniform) & --- & --- & seed reroll & human analyst \\
Readactor~\cite{crane2015readactor} & IR/machine & none & layout cost & by construction & layout randomization & memory disclosure \\
RandCompile~\cite{franzen2023randcompile} & IR & per-build random removal & --- & --- & build-to-build variance & forensic analysis \\
Wang et al.~\cite{bib4} & IR & RL sequence & 2-dim, in reward & --- & single sequence & known similarity detector \\
ReinObf~\cite{wei2025reinobf} & IR (block) & RL (PPO+GAT) & 2-dim, in reward & --- & single sequence & known detector ensemble \\
Polaris~\cite{zhou2026polaris} & MIR & none (fixed set) & --- & claimed & seed reroll & decompiler mid-end \\
AutoPass~\cite{li2026autopass} & IR & LLM agent & performance only & compiler oracle & --- (performance) & --- (performance) \\
\textbf{MIXCOMPILE (this paper)} & IR & \textbf{feature-driven, per-function} & \textbf{5-dim explicit} & \textbf{constructive + differential} & \textbf{combinatorial variant pool} & \textbf{amortizing / LLM-assisted} \\
\botrule
\end{tabular}
\end{table*}

%=======================================================
\section{MIXCOMPILE Method Design}
\label{sec:method}
%=======================================================

The vacancy identified in Section~\ref{sec:coordsys} defines the design target: a
compiler-level mechanism that selects obfuscation passes per program, under explicit
multi-dimensional cost control, with semantic safeguards and seed-independent
combinatorial diversity.

\subsection{Overview and Design Challenges}
\label{sec:overview}

Given an LLVM IR program, diversified compilation is the task of selecting a set of
obfuscation transformations for each program such that the resulting binaries behave
correctly while differing structurally under controlled runtime, size, and correctness
costs. Four challenges shape the design. First, uniform full activation maximizes
diversity but incurs unacceptable runtime overhead, code bloat, and inter-policy
interference, while static fixed passes provide no adaptation at all; the selection of
which combination fits each program must rest on measurable structural properties and
explicit cost estimates rather than one global setting. Second, applying passes
typically requires multiple IR traversals, which inflates compilation time and limits
applicability to large projects. Third, diversity gains must be weighed against
runtime, size, and correctness costs, and different deployment scenarios impose
different priorities, so the multidimensional trade-off must be encoded into a uniform
selection score while respecting user-specified deployment rules. Fourth, some passes
modify control flow and data flow in ways that can break program semantics if applied
carelessly, so the selection model must never enable them where correctness
constraints are not satisfied.

MIXCOMPILE addresses these challenges with a four-stage pipeline
(Figure~\ref{fig:pipeline}). An LLVM module enters the pipeline and passes
through four stages.

(1)~\textbf{Feature collection}. In a single traversal, all basic blocks
and instructions of each eligible function are visited; CFG structure counts,
instruction-type statistics, call flags, and inline-assembly flags are accumulated, and
an IBR def-use safety pre-screen is performed at the same time. After the traversal,
the average CFG depth is computed with a single BFS from the entry block, and
cyclomatic complexity plus the conditional-jump ratio are derived from the collected
counts. The program is never re-analyzed per pass.

(2)~\textbf{Applicability filtering}. Each function's (or basic block's)
feature tuple is checked against the structural preconditions of the seven passes.
Only passes whose minimal structural requirements are satisfied go on to the scoring
stage.

(3)~\textbf{Cost-aware selection}. Each candidate pass is scored
by the five-dimensional cost model under the configured deployment profile, with
budget gates and boundary randomization applied where appropriate
(Sections~\ref{sec:costmodel}--\ref{sec:boundary}).

(4)~\textbf{Correctness checking and annotation emission}. The candidates are ordered, IBR always first and the rest sorted by descending
score, then checked one by one against correctness constraints and written out as
LLVM Metadata annotations. Downstream passes read these annotations to decide whether
to apply their transformations; selection is thereby decoupled from execution.

MIXCOMPILE supports seven obfuscation strategies across two granularities. At the
function level: indirect branching (IBR), bogus control flow (BCF), control-flow
flattening (FLA), basic block splitting (SPLIT), and global variable indirection
(IGV). At the basic block level: instruction substitution (SUB) and indirect call
(ICALL). The user changes the sensitivity of the cost model by setting the deployment
profile (five weights) and the score threshold.

\paragraph{Granularity.}
Decisions are \emph{implemented} per function and per basic block, features are
collected at both levels and every unit is scored individually, and the per-function
pairing is empirically load-bearing: permuting which function carries which features
changes per-function decisions far beyond seed noise (Section~\ref{sec:e6},
experiment E6), and programs receive measurably different strategy unions
(Section~\ref{sec:rq3}, analysis T1).

\begin{figure}[!t]
\centering
\includegraphics[width=0.8\textwidth]{mixcompile_flow_revised.pdf}
\caption{Overall MIXCOMPILE pipeline: input LLVM module $\rightarrow$ four stages $\rightarrow$ Metadata-annotated IR module.}
\label{fig:pipeline}
\end{figure}

\subsection{Program Feature Extraction and Quantification}

MIXCOMPILE measures each function's CFG structure, branch density, instruction
types, and call behavior, and packs these measurements into a feature tuple. The
tuple is the only input used for building candidate pass sets and for cost
scoring. This makes decisions traceable: if a pass is turned on or off for a
particular function, we can point to which structural feature led to that
choice, and the full decision is recorded in an auditable log
(Section~\ref{sec:interface}).

\begin{definition}[Program Feature Collector]
Given an LLVM function $F$, MIXCOMPILE's feature collector produces a function-level feature
tuple $\Phi_F$ and a set of basic-block-level feature tuples $\{\varphi_B\}_{B\in F}$.
$\Phi_F$ characterizes the function's overall control-flow complexity, branch density,
instruction composition, call behavior, presence of inline assembly, and IBR compatibility;
$\varphi_B$ characterizes each basic block's local instruction scale, replaceable-instruction
ratio, and call presence.
\end{definition}

The function-level feature tuple is formally defined as:
\begin{equation}
\begin{aligned}
\Phi_F = \langle & \mathit{NumBBs}, \mathit{NumEdges}, \mathit{NumCondJumps}, \mathit{NumTotalJumps}, \\
                 & \mathit{CondJumpRatio}, \mathit{CyclomaticComplexity}, \mathit{AvgCFGDepth}, \\
                 & \mathit{TotalInsts}, \mathit{TotalSubInsts}, \mathit{MaxBBSize}, \mathit{SubRatio}, \\
                 & \mathit{HasCall}, \mathit{HasInlineAsm}, \mathit{IBRCompatible} \rangle
\end{aligned}
\label{eq:func_features}
\end{equation}

The basic-block-level feature tuple is formally defined as:
\begin{equation}
\varphi_B = \langle \mathit{BBSize}, \mathit{SubInsts}, \mathit{SubRatio}, \mathit{HasCall}, \mathit{HasInlineAsm} \rangle
\label{eq:bb_features}
\end{equation}

Table~\ref{tab:features} lists the function-level and basic-block-level features together
with their semantics and the passes they primarily influence.

\begin{table}[!t]
\centering
\footnotesize
\caption{Program feature summary}
\label{tab:features}
\begin{tabular*}{\textwidth}{@{\extracolsep\fill}lll@{}}
\toprule
Feature Field & Semantics & Primary Passes  \\
\midrule
NumBBs & Number of basic blocks in CFG & FLA, BCF  \\
NumEdges & Number of directed edges in CFG & FLA, BCF  \\
NumCondJumps & BBs ending with conditional branch & IBR \\
NumTotalJumps & BBs ending with any branch & IBR \\
CondJumpRatio & Conditional jump ratio  & IBR  \\
Cyclomatic Complexity & Cyclomatic complexity $M = E - N + 2$ & BCF \\
AvgCFGDepth & Mean BFS depth of all reachable BBs & FLA  \\
TotalInsts & Total IR instructions in function & ---  \\
TotalSubInsts & Total replaceable instructions in function & SUB \\
MaxBBSize & Max instruction count in any BB & SPLIT  \\
SubRatio & Replaceable instruction ratio & SUB \\
HasCall & Existence of CallInst or InvokeInst & IGV, ICALL \\
HasInlineAsm & Existence of inline assembly & IGV \\
IBR Compatible & IBR def-use safety check result & IBR \\
\botrule
\end{tabular*}
\end{table}

\subsubsection{Function-Level CFG Structural Features}

From the LLVM IR CFG, MIXCOMPILE pulls out four kinds of structural features: scale
features (NumBBs, NumEdges), branch-density features (NumCondJumps, NumTotalJumps,
CondJumpRatio), complexity features (CyclomaticComplexity), and hierarchical features
(AvgCFGDepth). All are computed directly from the CFG that LLVM already provides; no
extra data-flow analysis is needed. NumBBs and NumEdges indicate how large the CFG is,
and a bigger CFG generally leaves more room for several passes to do useful work.
CyclomaticComplexity measures the number of independent paths through the function and
is the main factor for deciding whether BCF makes sense: when the complexity is at
least 3, the function has at least two independent paths, so inserting bogus branches
can create effective confusion. CondJumpRatio characterizes conditional control-flow
density and is the core input for IBR decision making, because IBR produces a richer
distribution of jump targets in functions with dense conditional branches. AvgCFGDepth
is calculated from a BFS of the entry block over reachable basic blocks; deeply nested
functions flattened by FLA contrast more strongly with the original nested structure,
providing stronger anti-analysis effect.

\subsubsection{Basic-Block-Level and Instruction-Level Features}

As MIXCOMPILE walks through the instructions inside each basic block, it collects
the basic-block-level and instruction-level numbers that later feed SUB and ICALL
decisions. BBSize is the number of IR instructions in the block and is used for
scale constraints on both ICALL and SPLIT. For ICALL we only consider basic
blocks with fewer than 5 instructions, because the overhead of redirecting a call
is proportionally smaller in a short block. For SPLIT, the function needs at
least one basic block with 5 or more instructions; otherwise there is nothing
big enough to be worth splitting.

SubInsts counts how many instructions in the block are replaceable. Five
arithmetic/logic opcodes qualify: Add, Sub, And, Or, and Xor. These can be swapped out
for different instruction sequences that compute the same result. SUB can only be
applied if $\text{SubInsts} > 0$ in the block. SubRatio reflects the spatial density of
substitutions within a block and modulates the dynamic score of SUB: the higher the
ratio, the greater the local diversity gain.

HasCall is a flag for the presence of a CallInst or InvokeInst. For ICALL to be
considered, we need both $\text{HasCall} = \text{true}$ and $\text{BBSize} < 5$.
HasInlineAsm is a flag for inline assembly, aggregated up to the function level as a
safety filter for IGV: if a function has inline assembly anywhere, IGV is not applied
to it. The flag does not affect ICALL selection, since inline assembly matters for
IGV, not for ICALL. For call instructions the collector checks whether the called
function contains inline assembly; all these checks take constant time and do not
change the overall complexity of the traversal.

\subsubsection{IBR Compatibility Feature}

IBRCompatible is not a statistical feature but a Boolean feature checked by a
lightweight def-use construction. It shifts the correctness constraint of
control-flow-sensitive transformation into the feature extraction stage, so that the
subsequent cost model cannot select IBR on unsafe functions. The checks run
synchronously while traversing the use chain of instructions and apply three rejection
conditions:

(1)~\textbf{Same-block PHI use.} For each instruction $I$ that defines a value, we
look at all its uses. If any use is a PHI node in the same basic block as $I$, the
function is marked incompatible with IBR.

(2)~\textbf{Cross-block use in a multi-predecessor block.} If any use of $I$ is in a
different basic block, and that block has more than one predecessor, the function is
rejected for IBR.

(3)~\textbf{Conditional branch on the unique-predecessor chain.} If the use is in a
block $U_B$ that has exactly one predecessor $P$, we trace backward along this chain
of single-predecessor blocks until we reach the block $B$ where $I$ was defined. If
any block along the chain ends with a conditional branch, or has more than one
predecessor, IBR is rejected.

Together, these checks restrict IBR to functions where def-use relationships have a
clean, linear structure, so that replacing direct branches with indirect ones cannot
compromise the program's semantics.

\subsubsection{Feature Extraction Algorithm}

The feature collection is done in a single pass through each function, light enough
that it does not noticeably slow down compilation. Algorithm~\ref{alg:feature} spells
out the full procedure.

\begin{algorithm}[!t]
\footnotesize
\caption{MIXCOMPILE Feature Collection (FeatureCollection)}
\label{alg:feature}
\begin{algorithmic}[1]
    \Require LLVM function $F$
    \Ensure Function-level feature tuple $\Phi_F$, basic-block-level feature map $\{\varphi_B\}$
    \State Initialize $\Phi_F$, all fields to 0 or false
    \State $\text{IBRCompat} \gets \text{true}$
    \For{each basic block $B \in F$}
        \State $\Phi_F.\text{NumBBs} \gets \Phi_F.\text{NumBBs} + 1$
        \State $\Phi_F.\text{NumEdges} \gets \Phi_F.\text{NumEdges} + \text{SuccCount}(B)$
        \State Update jump statistics from $\text{Terminator}(B)$:
        \If{is conditional branch} \State $\Phi_F.\text{NumCondJumps} \gets \Phi_F.\text{NumCondJumps} + 1$ \EndIf
        \If{is any branch} \State $\Phi_F.\text{NumTotalJumps} \gets \Phi_F.\text{NumTotalJumps} + 1$ \EndIf
        \State Initialize local feature $\varphi_B$, all fields to 0 or false
        \For{each instruction $I \in B$}
            \State $\varphi_B.\text{BBSize} \gets \varphi_B.\text{BBSize} + 1$
            \If{$\text{IsSubstitutable}(I)$} \State $\varphi_B.\text{SubInsts} \gets \varphi_B.\text{SubInsts} + 1$ \EndIf
            \If{$I$ is CallInst or InvokeInst}
                \State $\varphi_B.\text{HasCall} \gets \text{true}$
                \If{$I$ contains inline assembly} \State $\varphi_B.\text{HasInlineAsm} \gets \text{true}$ \EndIf
            \EndIf
            \If{$I$ has uses and $\text{IBRCompat} = \text{true}$}
                \State $\text{IBRCompat} \gets \text{IBRCompat} \land \text{CheckIBRUseSafety}(I, B)$
            \EndIf
        \EndFor
        \State $\varphi_B.\text{SubRatio} \gets \text{SafeRatio}(\varphi_B.\text{SubInsts}, \varphi_B.\text{BBSize})$
        \State Accumulate to function-level totals and update $\Phi_F.\text{MaxBBSize}$,
               $\Phi_F.\text{HasCall}$, $\Phi_F.\text{HasInlineAsm}$
    \EndFor
    \State $\Phi_F.\text{CyclomaticComplexity} \gets \max(\text{NumEdges} - \text{NumBBs} + 2, 1)$
    \State $\Phi_F.\text{CondJumpRatio} \gets \text{SafeRatio}(\text{NumCondJumps}, \text{NumTotalJumps})$
    \State $\Phi_F.\text{AvgCFGDepth} \gets \text{BFSAverageDepth}(F)$
    \State $\Phi_F.\text{SubRatio} \gets \text{SafeRatio}(\text{TotalSubInsts}, \text{TotalInsts})$
    \State $\Phi_F.\text{IBRCompatible} \gets \text{IBRCompat}$
    \Return $\Phi_F, \{\varphi_B\}$
\end{algorithmic}
\end{algorithm}

For each function $F$, MIXCOMPILE completes all feature extraction within a single
traversal. After traversal, cyclomatic complexity is computed directly as
$\text{NumEdges} - \text{NumBBs} + 2$, the conditional-jump ratio via
$\text{SafeRatio}(\text{NumCondJumps}, \text{NumTotalJumps})$, and the average CFG
depth via a single BFS. The per-function feature-extraction cost is $O(|I| + |E| +
|U|)$, where $|I|$ is the number of IR instructions, $|E|$ the number of CFG edges,
and $|U|$ the number of def-use relationships.

\subsection{Cost-Aware Selection Model}
\label{sec:costmodel}

After feature extraction, MIXCOMPILE estimates the benefit and cost of each applicable
pass for the current program structure, then reduces the five dimensions to a scalar
score under a selected deployment profile. The separation matters: the JSON file
stores the pass-specific base vectors, whereas the profile specifies the
deployment-level trade-off among security, diversity, runtime, size, and risk.

\subsubsection{Pass Base Cost Vector}

Each pass $P$ is associated with a fixed five-dimensional base cost vector
$C_P^{\text{base}}$ encoding its intrinsic properties:
\begin{equation}
C_P^{\text{base}} = \langle s_P, d_P, r_P, z_P, k_P \rangle ,
\label{eq:cost_vector}
\end{equation}
where $s_P$ and $d_P$ are security and diversity rewards, and $r_P$, $z_P$, and
$k_P$ are runtime cost, size cost, and correctness risk. These values are base
quantities; the feature-modulated terms of Section~3.3.3 turn them into
per-function scores. Table~\ref{tab:optimized_pass_vectors}
reports the vectors frozen for the formal optimized evaluation.

\begin{table*}[!t]
\centering
\footnotesize
\caption{Frozen optimized pass base vectors.}
\label{tab:optimized_pass_vectors}
\begin{tabular}{lccccc}
\toprule
Pass & Security gain & Diversity gain & Runtime cost & Size cost & Correctness risk \\
\midrule
SUB   & 0.490126 & 0.199157 & 0.075497 & 0.048800 & 0.108150 \\
BCF   & 0.434216 & 0.871526 & 0.773944 & 0.709331 & 0.254053 \\
FLA   & 0.707056 & 0.488146 & 0.736028 & 0.215843 & 0.365736 \\
SPLIT & 0.306698 & 0.172056 & 0.108178 & 0.094880 & 0.114597 \\
IBR   & 0.572982 & 0.597764 & 0.242986 & 0.236418 & 0.520390 \\
ICALL & 0.458715 & 0.157404 & 0.169846 & 0.106568 & 0.376419 \\
IGV   & 0.516724 & 0.362678 & 0.104081 & 0.135240 & 0.530250 \\
\bottomrule
\end{tabular}
\end{table*}

\subsubsection{Weight Configuration}

The deployment profile is
\begin{equation}
W = \langle w_{\mathrm{sec}}, w_{\mathrm{div}}, w_{\mathrm{run}},
w_{\mathrm{size}}, w_{\mathrm{risk}} \rangle .
\label{eq:weights}
\end{equation}
The reported experiments use the optimized deployment profile
\begin{equation}
W_{\mathrm{opt}}=\langle1.10,1.05,0.90,0.25,0.95\rangle .
\label{eq:optimized_weights}
\end{equation}
The profile values are calibrated engineering parameters; Section~4 reports the
calibration, sensitivity, and ablation evidence behind them.

\subsubsection{Feature-Modulated Dynamic Scoring}

Although a selected profile is fixed for one experiment, MIXCOMPILE's score is
dynamic. For a candidate pass $P$ and feature tuple $x$ (where $x$ is either $\Phi_F$
or $\varphi_B$), the composite score works as follows:
\begin{equation}
\text{Score}(P, x, W) = w_{\mathrm{sec}} \cdot S_P(x) + w_{\mathrm{div}} \cdot D_P(x) - w_{\mathrm{run}} \cdot R_P(x) - w_{\mathrm{size}} \cdot Z_P(x) - w_{\mathrm{risk}} \cdot K_P(x)
\label{eq:score}
\end{equation}

The terms are dynamic values modulated according to the corresponding features, so the
same pass can produce completely different scores on different characteristics. The
specific modulation rules are as follows:

\textbf{IBR:} $S_{\text{IBR}}$ increases with CondJumpRatio, as denser conditional
branching yields richer jump-target diversity; $R_{\text{IBR}}$ likewise increases
with conditional-jump ratio as more branch points must be indirected.

\textbf{BCF:} BCF uses a workload-aware estimate.  Let $N_{BB}$ be the number of basic
blocks, $N_I$ the number of instructions, $c$ the cyclomatic complexity, $q$ the configured
BCF probability in percent, and $L$ the configured loop count.  We compute
\begin{align}
f_{\mathrm{BCF}} &= \operatorname{clamp}_{[0,1]}\left(\frac{c-2}{4}\right),
&p &= \operatorname{clamp}_{[0,1]}(q/100), \\
E_{\mathrm{sum}} &= \sum_{i=0}^{L-1}(1+3p)^i,
&E_{BB} &= N_{BB}\,p\,E_{\mathrm{sum}}, \label{eq:bcf_workload}\\
m_R &= 1+\frac{\log_2(1+18E_{BB})}{10},
&m_Z &= 1+\frac{\log_2(1+E_{BB}(N_I/N_{BB}+18))}{12}.
\label{eq:bcf_multipliers}
\end{align}
The adjusted BCF dimensions are
\begin{align}
S_{\mathrm{BCF}}&=s_{\mathrm{BCF}}f_{\mathrm{BCF}}, &
D_{\mathrm{BCF}}&=d_{\mathrm{BCF}}f_{\mathrm{BCF}}, \\
R_{\mathrm{BCF}}&=r_{\mathrm{BCF}}f_{\mathrm{BCF}}m_R, &
Z_{\mathrm{BCF}}&=z_{\mathrm{BCF}}f_{\mathrm{BCF}}(1+\mathrm{CondJumpRatio})m_Z, \\
K_{\mathrm{BCF}}&=k_{\mathrm{BCF}}f_{\mathrm{BCF}}. &&
\end{align}
Thus BCF cost reflects its configured probability, repeated expansion, original block
size, and expected added instructions, instead of relying only on a static pass
vector. BCF is eligible only when $c\geq3$, $N_I\leq2000$, and $E_{BB}\leq500$.  The
latter two conditions are explicit budget gates that limit long-tail compile-time,
runtime, and size expansion.

\textbf{FLA:} $S_{\text{FLA}}$ increases with AvgCFGDepth, as flattening produces
more thorough structural disruption in deeply nested functions;
$R_{\text{FLA}}$ increases with cyclomatic complexity;
$Z_{\text{FLA}}$ increases with the number of basic blocks.

\textbf{SPLIT:} $S_{\text{SPLIT}}$ increases with MaxBBSize; $Z_{\text{SPLIT}}$
increases with maximum basic-block size.

\textbf{SUB:} $S_{\text{SUB}}$ increases with the basic block's SubRatio; runtime and
size costs increase with substitution ratio and count.

\textbf{ICALL:} Generates benefits and costs only on basic blocks where
$\text{HasCall} = \text{true}$ and $\text{BBSize} < 5$.

\textbf{IGV:} When $\text{HasInlineAsm} = \text{true}$, correctness risk is set to a
high value, effectively equivalent to direct exclusion by the applicability predicate.
IGV is always placed in the boundary state, because the available static features do
not distinguish strong from weak IGV candidates.

After feature-modulated dynamic scoring, MIXCOMPILE applies a global threshold $\tau$
as a final filter: only candidates with $\text{Score} > \tau$ enter the candidate set.
$\tau$ is exposed via the \texttt{--mix-score-threshold} command-line option with a
default value of 0.0.

\subsection{Controlled Randomization via Boundary Conditions}
\label{sec:boundary}

Deterministic feature rules ensure the interpretability of decisions, but if all
boundary samples make the same decision, repeated compilation of the same program may
still produce highly similar results, re-enabling attack amortization. MIXCOMPILE
therefore applies a random admission mechanism in the boundary region where a pass is
applicable but its benefit is uncertain.

\begin{definition}[Boundary Predicate]
For a pass $P$, $\text{Boundary}(P, x)$ is true if and only if the dominant feature of
$P$ lies within a predefined boundary interval close to its applicability threshold.
\end{definition}

The specific boundary intervals are listed in Table~\ref{tab:boundary}.

\begin{table}[!t]
\centering
\footnotesize
\caption{Pass boundary intervals}
\label{tab:boundary}
\begin{tabular}{ccc}
\toprule
Pass & Boundary Interval & Condition Description \\
\midrule
IBR & $0.15 \leq \text{CondJumpRatio} < 0.25$ & Benefit uncertain at low conditional-branch density \\
BCF & $3 \leq \text{CyclomaticComplexity} < 5$ & Bogus-branch differentiation insufficient at low complexity \\
FLA & $2.0 \leq \text{AvgCFGDepth} < 3.0$ & Structural change limited at shallow CFG depth \\
SPLIT & $5 \leq \text{MaxBBSize} < 8$ & Splitting benefit limited for smaller basic blocks \\
SUB & $\text{SubRatio} < 0.5 \lor \text{SubInsts} < 3$ & Replaceable-instruction density or count insufficient \\
IGV & --- & Always placed in boundary state after applicability filtering \\
ICALL & --- & Always placed in boundary state after applicability filtering \\
\botrule
\end{tabular}
\end{table}

Boundary randomization triggers only when three conditions hold at the same time: the
pass satisfies all structural preconditions, its net benefit exceeds the threshold, and
its dominant feature lies in the boundary interval. Only then does MIXCOMPILE call a
lightweight randomized decision function to decide whether the pass is still included.

$\text{ChooseMachine}()$ is an approximately 50\% binary random back-off; the exact
probability is not a tuning parameter. When the same source code is compiled
multiple times, this mechanism means that different compilation runs obtain different
pass combinations, producing binaries that are structurally diverse. The scope of
randomization is strictly limited to the boundary region where the feature signal is
weak, so interpretability is preserved where features carry information and diversity
is injected where they do not.

The design also admits an information-theoretic reading~\cite{shannon1948mathematical}.
Whether a boundary-region pass is applied is a binary random variable $X$ with
$\Pr(X{=}1) \approx 0.5$, so each boundary decision contributes close to the maximum
one bit of entropy, $H(X) \approx 1$, and the entropy is maximized precisely where
the feature signal is weakest. An attacker observing a deployed binary therefore
faces maximum uncertainty about the compilation decisions in exactly the region where
deterministic inference from program structure fails. Outside the boundary region,
features carry the information and the decision is deterministic; inside it, the
entropy of the decision cannot be reduced without the compilation seed, which the
threat model withholds from the adversary. Diversity is thus injected at the point of
maximum ignorance rather than uniformly.

\subsection{Correctness Protection}

Applying passes to the same code region can lead to compilation errors. To avoid
this, MIXCOMPILE applies dedicated correctness checks, described here together with
the remaining safeguards.

\subsubsection{IBR Def-Use Safety Analysis}

The IBR transformation replaces a direct branch with an indirectly computed target.
If a subsequent transformation changes the control flow around the branch, def-use
chains may no longer be dominated, and semantics can break. To prevent this, MIXCOMPILE
performs the structural def-use safety analysis during the feature extraction phase
(Algorithm~\ref{alg:ibr_check}), so that the selection model can never enable IBR on a
function where the analysis fails.

\begin{algorithm}[!t]
\footnotesize
\caption{IBR Compatibility Check (CheckIBRCompatibility)}
\label{alg:ibr_check}
\begin{algorithmic}[1]
    \Require Function $F$
    \Ensure true if IBR can be safely applied, false otherwise
    \For{each basic block $B \in F$}
        \For{each instruction $I \in B$}
            \For{each use $U$ of $I$}
                \If{$U$ is a PHI node and $\text{Parent}(U) = B$}
                    \Return false
                \EndIf
                \If{$U$ is an instruction in another block $U_B$}
                    \If{$\text{PredCount}(U_B) > 1$} \Return false \EndIf
                    \State $P \gets \text{UniquePredecessor}(U_B)$
                    \If{$\neg \text{CheckPredChain}(P, B)$} \Return false \EndIf
                \EndIf
            \EndFor
        \EndFor
    \EndFor
    \Return true
\end{algorithmic}
\end{algorithm}

The algorithm checks the control flow of each function, performing CheckPredChain on
variables whose use site and def site are in different basic blocks. CheckPredChain
requires the existence of a unique, unconditionally branched predecessor chain from
the use block back to the definition block. This constraint guarantees that IBR is not
enabled in def-use regions that rely on multipath confluence or conditional control
flow, which are the structural patterns where indirect branches are most likely to
introduce semantic errors.

\subsubsection{Conflict Resolution and Additional Safeguards}

During candidate ranking, IBR is given priority. In particular, an IBR--FLA conflict
is resolved by retaining IBR and removing FLA.  If both BCF and FLA survive their
initial eligibility and threshold checks and no IBR--FLA conflict exists, the
implementation re-estimates the two control-flow scores before retaining one of them.
This second comparison removes the initial clamps:
\begin{equation}
f^{\prime}_{\mathrm{BCF}}=(c-2)/4,\qquad
f^{\prime}_{\mathrm{FLA}}=\mathrm{AvgCFGDepth}/3,
\end{equation}
and the FLA runtime factor becomes $1+c/6$.  The five scaled dimensions are also left
unclamped in this comparison.  This preserves structural differences among highly
complex functions that would otherwise collapse to the same value in $[0,1]$.  The
pass with the greater comparison score is retained; an exact tie is resolved using the
initial score.  Remaining candidates are processed in descending score order while
respecting the conflict relation.

In addition to IBR validation, MIXCOMPILE enforces two safeguards. First, IGV must
not be used on any function that contains inline assembly: global variable indirection
rewrites variable references and does not preserve the implicit memory-access
conventions and register-allocation assumptions of handwritten assembly. Second, a
SPLIT--ICALL conflict is resolved by one-way suppression: when SPLIT has been included
in the function-level pass set, ICALL candidates for all basic blocks within that
function are automatically excluded, because basic block splitting may eliminate call
instructions targeted by ICALL or change basic block boundaries.

\paragraph{Semantic guarantee.}
The guarantees above are constructive: for well-defined programs, the selected
transformations preserve semantics by construction, and every formal run
additionally passes differential verification against the unobfuscated build
(Section~\ref{sec:correctness}).

\subsection{LLVM Metadata-Based Pass Annotation}

MIXCOMPILE decouples policy selection from policy execution. The PassDecider component
is a policy selector and performs no code transformation; downstream passes are
transformation executors that decide whether to run based on annotations. This design
follows the single-responsibility principle and lets new passes plug in without
touching the selection core: a new pass declares its base weights, applicability
predicate, and conflict rules, and neither the core logic nor any existing pass needs
modification.

At the function level, MIXCOMPILE emits five annotation kinds in the form of LLVM
function metadata, corresponding to IBR, IGV, FLA, BCF, and SPLIT, each supporting
positive and negative annotation. At the basic block level, MIXCOMPILE emits two
annotation kinds, SUB and ICALL, in the form of LLVM basic-block metadata, written only
as positive annotations when the pass is selected.

Downstream passes query the corresponding annotation through LLVM's metadata API at
their entry point: if a positive annotation is present, the transformation is
performed; otherwise the pass skips the function or block. Annotation reading is an
$O(1)$ string comparison and does not increase the runtime overhead of the passes
themselves.

Annotation-based design offers two further engineering advantages. First,
testability: the selection logic can be unit-tested by inspecting annotation
output without executing a full compilation pipeline. Second, efficiency:
downstream passes do not need to traverse every function and block to re-derive
decisions but apply policy directly according to annotations, reducing
compilation time.

\subsection{Machine-Invocable Strategy Interface}
\label{sec:interface}

Because selection is decoupled from execution and every decision is logged, the
strategy space is exposed in a form that external callers can consume. Three interface
elements are relevant.

\textbf{(1) Restricted strategy space.} The pass base vectors
(Table~\ref{tab:optimized_pass_vectors}), the deployment profile
(Eq.~\eqref{eq:weights}), and the threshold $\tau$ together define a finite,
well-documented space of legal pass combinations. An external caller, whether a
script, a search procedure, or an LLM-based compilation agent, can enumerate and
reason about this space without reverse-engineering the compiler.

\textbf{(2) Interpretable compilation feedback.} Every function-level control-flow
decision emits a \texttt{[MIXCOMPILE][CF\_DECISION]} record containing the function
name, block and instruction counts, cyclomatic complexity, conditional-jump ratio,
mean CFG depth, BCF probability and loop count, $E_{BB}$, both BCF budget limits,
budget eligibility, runtime and size multipliers, the initial BCF/FLA totals, the
unclamped comparison totals, and the selected pass. Budget rejections and BCF--FLA
arbitration are therefore auditable after the fact.

\textbf{(3) Stable invocation surface.} The profile is supplied as a JSON file, the
threshold via \texttt{--mix-score-threshold}, and the randomization seed via
\texttt{-mix-seed}; the combination (seed, profile, JSON vectors, BCF budgets, target
binary and optimization level) fully determines the output and constitutes the
reproduction contract of Section~\ref{sec:setup}.

On this interface, an LLM agent plans pass configurations in the structured
strategy space, invokes the compiler, and verifies the outcome against the
per-decision feedback; the E10 experiment (Section~\ref{sec:ongoing}) shows that
this usage preserves semantics in 98.9\% of cases, versus 63.2\% for direct LLM
source rewriting under identical verification.

\subsection{Algorithm Workflow and Complexity Analysis}

The MIXCOMPILE main workflow is given in Algorithm~\ref{alg:main}.

\begin{algorithm}[!t]
\footnotesize
\caption{MIXCOMPILE Main Workflow}
\label{alg:main}
\begin{algorithmic}[1]
    \Require LLVM module $M$, threshold $\tau$
    \Ensure Annotated LLVM module
    \State $W \gets \text{SelectedProfile}()$
    \For{each function $F \in M$}
        \If{$\neg \text{Eligible}(F)$} \State \textbf{continue} \EndIf
        \State $\text{SplitCriticalEdges}(F)$
        \State $\Phi_F, \{\varphi_B\} \gets \text{CollectFeatures}(F)$
        \State $\mathcal{S}_F \gets \text{SelectFunctionPasses}(\Phi_F, W, \tau)$
        \State $\text{AnnotateFunction}(F, \mathcal{S}_F)$
        \For{each basic block $B \in F$}
            \State $\mathcal{S}_B \gets \text{SelectBasicBlockPasses}(\varphi_B, W, \tau, \mathcal{S}_F)$
            \State $\text{AnnotateBasicBlock}(B, \mathcal{S}_B)$
        \EndFor
    \EndFor
    \Return $M$
\end{algorithmic}
\end{algorithm}

We look at the computational cost of MIXCOMPILE at two levels.
\textbf{Per-function level:} feature collection does one traversal over all basic
blocks and instructions, costing $O(|I| + |E| + |U|)$ where $|I|$ is instructions,
$|E|$ is CFG edges, and $|U|$ is def-use chains. Scoring computes one weighted sum
for each of the $|P| = 7$ passes, which is $O(1)$. Sorting the candidates and
resolving conflicts also takes $O(1)$ since there are at most 7 candidates.
Basic-block-level selection repeats scoring and boundary checking for at most 2
passes per basic block, also constant work. Writing the annotations is a constant
number of metadata strings per function and basic block.
\textbf{Whole-module level:} adding up over all functions, the total cost grows
linearly with the size of the IR. MIXCOMPILE requires no separate IR traversal
per pass: all seven passes' feature needs are covered in a single traversal,
which keeps compilation time linear and predictable.

%=======================================================
\section{Experimental Evaluation}
%=======================================================

This section answers five research questions.

\subsection{Research Questions}

\textbf{RQ1 (Model behavior):} Do the five cost dimensions carry distinct roles, and
are the frozen weights defensible rather than arbitrary?

\textbf{RQ2 (Where features act):} Do program features carry decision-relevant
information, and at what granularity, per-function pairing or marginal distribution?

\textbf{RQ3 (Diversity shape):} Are repeated compilations of the same program mutually
distinct, and do different programs receive different strategies?

\textbf{RQ4 (Deployment cost):} What runtime cost does the optimized profile incur on
real AI-inference and cryptographic workloads, relative to GCC, LLVM, and full
obfuscation?

\textbf{RQ5 (Analysis cost):} Does the output measurably increase the cost of
automated binary analysis?

Two pre-registered LLM-era studies (E9, E10) complement these questions from
the offense and defense sides: E9 measures what an LLM attacker can actually
do to the compiled output, and E10 measures what an LLM planner can safely do
through the strategy interface.

\subsection{Experimental Setup}
\label{sec:setup}

The mechanism experiments (E1--E7, diversity analyses) run on a frozen program split
of 44 tuning / 43 validation / 43 test programs (program-level disjoint,
\texttt{split\_seed=20260903}), drawn from the LLVM test-suite and benchmark
families; weight selection was pre-registered on the validation split and
confirmed once on the test split. Correctness verification compares obfuscated
builds against unobfuscated builds under the benchmark test inputs. The workload
and analysis-cost measurements (llama.cpp, OpenSSL, cross-compiler radiff2, Ghidra
decompilation time) use the optimized pass-vector configuration with seed 1,
$\tau=0.0$, and the BCF gates $N_I\leq2000$ and $E_{BB}\leq500$, and passed output
validation on all 31 binaries of the LLVM test-suite validation set.

All experiments run on a workstation with an Intel Core i9-12900H (x86\_64, 14
physical cores, up to 5.0\,GHz), 64\,GB RAM, and 64-bit Ubuntu Linux 22.04. Compared
compilers: GCC~11.4.0, LLVM/Clang~17.0.6, OLLVM (built on LLVM~17 with all available
passes enabled), and MIXCOMPILE (built on LLVM~17.0.6, PassDecider handling
selection).

Every formal binary is fully determined by the seed, profile JSON, pass vector
JSON, BCF budget gates, and target binary/optimization level, and every
selection decision is recorded in a \texttt{CF\_DECISION} log archived with
the binaries; recompiling with identical inputs reproduces the decision logs
and the output binaries exactly. Repeated throughput measurements (llama.cpp,
OpenSSL) are summarized with the arithmetic mean; heterogeneous, right-skewed
quantities (radiff2 similarities, decompilation-time ratios) are summarized
with medians. Within one comparison, the same statistic is applied to every
compiler.

\subsection{RQ1: Cost-Model Behavior (E1--E5)}

\subsubsection{Weight benchmark against random weights (E1)}
The five weights are engineering parameters; the question is whether they carry
information. We compared the pre-registered optimized profile against 30 random legal
weight configurations (plus fixed baselines) over 43 programs $\times$ 30 seeds,
scoring every configuration on all five cost dimensions (Table~\ref{tab:e1_weights}).

\begin{table}[!t]
\centering
\footnotesize
\caption{E1: weight benchmark. Random-weight configurations fall into four behavior
bands; the pre-registered optimized profile is shown for comparison.}
\label{tab:e1_weights}
\begin{tabular}{lrrll}
\toprule
Configuration & Count & Activations & Exec (median) & Correctness \\
\midrule
Random: not activated & 4 & 0 & $\approx$1.0$\times$ & 1.0 \\
Random: low activation & 4 & 6 & --- & 1.0 \\
Random: mid band & 18 & 10--24 & 1.003--1.046$\times$ & see note \\
Random: high activation & 2 & 28--29 & $+$1.2--1.7\% & 1.0 \\
\textbf{optimized (pre-reg.)} & --- & \textbf{27.0} & \textbf{1.0444$\times$} & \textbf{1.0} \\
\botrule
\end{tabular}
\end{table}

Three findings. \textbf{(1)~Dominance.} The optimized profile dominates 28 of the 30
random candidates on the pre-registered validation criteria (test-split confirmation:
correctness 1.0, median exec 1.0444$\times$, median activation 27.0); the two
undominated candidates (random\_015/016) achieve higher activation at smaller code
size, so the pre-registered selection occupies an endpoint of the Pareto front.
\textbf{(2)~Random weights are risky.} Four configurations (random\_008/009/019/029)
silently corrupt correctness (pass rate 0.986--0.992; these are also the source of the
53 compiler-internal-error observations in the \texttt{voronoi} program); the cost
model's value is that it makes this trade-off explicit before compilation.
\textbf{(3)~The four bands} show the behavior space of unconstrained weights, from
complete inactivity (score threshold never met) to near-maximal activation with the
corresponding size bill ($1.18\times$). An equal-weight profile performs comparably
on some axes but was not the pre-registered selection.

\subsubsection{Dimension roles (E2)}
Removing each weight dimension in turn (on the validation split) reveals distinct,
non-redundant roles (Table~\ref{tab:e2_roles}).

\begin{table}[!t]
\centering
\footnotesize
\caption{E2: effect of removing each cost dimension.}
\label{tab:e2_roles}
\begin{tabular}{llll}
\toprule
Dimension removed & Activations & Exec (median) & Size (median) \\
\midrule
security & $-9$; 6 programs $\to$ 0 & --- & --- \\
runtime & --- & 1.07$\times$ & 1.25$\times$ \\
size & --- & 1.13$\times$ & 1.70$\times$ \\
risk & --- & 1.08$\times$ & 1.49$\times$ \\
diversity & $+1$ & $\approx$ unchanged & $\approx$ unchanged \\
\botrule
\end{tabular}
\end{table}

Security is the accelerator: without it, most selections collapse to zero. Runtime,
size, and risk are brakes with different strengths; each removal inflates both cost
axes, with the size dimension the strongest brake against code bloat. The diversity
dimension has the smallest marginal effect ($+1$ activation). The five dimensions
are not interchangeable: none of the four brakes can substitute for another
without shifting the cost profile.

\subsubsection{Per-pass calibration (E4)}
The pass base vectors are calibrated, not invented: each pass was enabled alone, and
its measured cost profile feeds the model coefficients (Table~\ref{tab:e4_perpass}).

\begin{table}[!t]
\centering
\footnotesize
\caption{E4: single-pass calibration. Units are the decision targets logged by the
calibration pipeline (functions or basic blocks per pass); ratios are medians relative
to the LLVM baseline.}
\label{tab:e4_perpass}
\begin{tabular}{lrrl}
\toprule
Pass & Eligible units & Activation rate & Exec ratio (median) \\
\midrule
SUB & 7362 & 81.1\% & 0.8065$\times$ \\
BCF & 2060 & 78.3\% & 1.0415$\times$ \\
FLA & 208 & 17.8\% & 1.0265$\times$ \\
SPLIT & 5164 & 100\% & 1.0241$\times$ \\
IBR & 480 & 30.6\% & 1.0203$\times$ \\
ICALL & 2528 & 65.0\% & 0.7792$\times$ \\
IGV & 2798 & 74.0\% & 1.0152$\times$ \\
\botrule
\end{tabular}
\end{table}

Two annotations. FLA activates rarely under single-pass conditions (17.8\%), so
its base vector rests on fewer observations than the other passes. SUB and
ICALL show median exec ratios \emph{below} 1.0; these measurements are
reported as recorded.

\subsubsection{Stability and interactions (E3, E5)}
Sixty-one perturbed weight configurations (single-dimension $\pm10\%$ to $\pm30\%$ and
30 joint perturbations, validation split) reproduce the original strategy with
consistency 0.888 (single-dimension) and 0.832 (joint) measured as per-program
pass-set agreement: the decisions are stable in the neighborhood of the frozen
weights. Four representative pass pairs
(FLA$\times$BCF, FLA$\times$SPLIT, IBR$\times$IGV, SPLIT$\times$ICALL; 440
measurement units each) show no super-additive effect, with combined outcomes
deviating from the sum of individual effects by $-0.8\%$ to $-2.6\%$; per-program
selection therefore rests on the cost asymmetry demonstrated in E2/E4.

\subsection{RQ2: Where Features Act (E6)}
\label{sec:e6}

E6 tests whether program features actually drive selection. Four conditions were
compared on the validation split: \emph{complete} (features + costs, the deployed
pipeline); \emph{feature-only} (cost vectors flattened); \emph{cost-only} (features
flattened); and \emph{shuffled-feature} (the function-to-feature assignment is
randomly permuted, with feature marginals, weights, seeds, and pass pool unchanged,
$\geq$30 permutations).

The evidence is the paired per-function Jaccard between each ablated condition and
the complete pipeline (Table~\ref{tab:e6_conditions}, per-function column).
Permuting which function carries which features moves the Jaccard to 0.549
($n=15{,}540$ function-level comparisons), well below the same-configuration
cross-seed floor of 0.708 ($n=225{,}330$): the function--feature pairing, not the
seed, determines which passes a function receives. Cost-only flattening diverges
to a similar degree (0.551), while feature-only flattening stays within seed
noise (0.724): the per-function placement is carried by the features and
modulated by the costs. The outcome columns separate the conditions as designed,
with cost-only cheapest on both axes. Figure~\ref{fig:e6_granularity} shows the
contrast graphically.

\begin{table}[!t]
\centering
\footnotesize
\caption{E6: four-condition ablation. $J_{\mathrm{func}}$: paired per-function
Jaccard against the complete condition; the same-configuration cross-seed floor
at this granularity is 0.708.}
\label{tab:e6_conditions}
\begin{tabular}{lrrr}
\toprule
Condition & Exec & Size & $J_{\mathrm{func}}$ \\
\midrule
complete & 1.217$\times$ & 1.487$\times$ & (0.708 cross-seed floor) \\
feature\_only & 1.352$\times$ & 1.252$\times$ & 0.724 \\
cost\_only & 1.061$\times$ & 1.000$\times$ & 0.551 \\
shuffled & 1.163$\times$ & 1.234$\times$ & 0.549 \\
\botrule
\end{tabular}
\end{table}

\begin{figure}[!t]
\centering
\includegraphics[width=0.9\textwidth]{e6_granularity.pdf}
\caption{E6 paired per-function Jaccard against the complete condition: shuffled
(0.549) and cost-only (0.551) fall below the same-configuration cross-seed floor
(dashed line, 0.708), while feature-only (0.724) stays within seed noise. Boxes:
interquartile range; white diamonds: means; whiskers: most extreme values within
1.5$\times$IQR.}
\label{fig:e6_granularity}
\end{figure}

Three conclusions follow. (1)~The pipeline is feature-driven at the function
level: both feature-ablating conditions fall below the seed floor, while
feature-only stays within it. (2)~Per-function targeted assignment is
mechanically real: the specific pairing of a function with its own features
determines which passes it receives. (3)~A security-side measurement of
targeted versus untargeted placement at matched overhead is the immediate
follow-up experiment.

\subsection{RQ3: Diversity Shape (E7, T1, T2)}
\label{sec:rq3}

\subsubsection{Variant-level uniqueness (E7)}
For each of the 12 security-subset programs, 30 optimized variants (seeds 1--30) were
compiled. Every program's 30 variants are pairwise distinct at the \texttt{.text}
level: 12/12 programs, 100\% uniqueness. Structural distance to the unobfuscated build
(radiff2, per-program median across variants): range 0.154--0.912, median 0.39, mean
0.45. The dispersion is itself informative: the least-affected program (12-IOtest,
0.912) and the most-affected (dc, 0.154) differ by roughly $6\times$, and the distance
does not correlate obviously with security sensitivity (security-sha sits at 0.680,
toward the ``similar'' side). Opcode 3-gram Jaccard against the unobfuscated build
ranges 0.548--1.000 (mean 0.76). Diversity is real but program-dependent, and we
report the spread rather than only the median. Table~\ref{tab:e7_diversity}
summarizes the RQ3 metrics.

\begin{table}[!t]
\centering
\footnotesize
\caption{RQ3 diversity summary (optimized arm).}
\label{tab:e7_diversity}
\begin{tabular}{ll}
\toprule
Metric & Result \\
\midrule
\texttt{.text} uniqueness (30 variants per program) & 12/12 programs, 100\% \\
radiff2 variant-to-unprotected, per-program median & 0.39 (range 0.154--0.912) \\
opcode 3-gram Jaccard vs.\ unprotected & mean 0.76 (range 0.548--1.000) \\
Between- vs.\ within-program strategy distance (T1) & $3.3\times$ (0.442 vs.\ 0.134) \\
Per-program cross-seed strategy uniqueness (T2) & 0.058 (stable within a program) \\
\bottomrule
\end{tabular}
\end{table}

\subsubsection{Inter-program strategy dispersion (T1)}
The strategies, that is, the pass combinations selected per program, differ
across programs. The distance between programs' modal pass-set unions is 0.442,
against a within-program cross-seed distance of 0.134: programs receive
significantly different strategies, at $3.3\times$ the seed-noise dispersion.

The complement is T2: within a program, cross-seed union uniqueness is 0.058, so a
program's strategy is stable across reseeds. Read together, the strategy is a stable,
program-specific fingerprint: selection is driven by the program's features rather
than by the seed. The strategy repertoire is coarse (7--15 distinct unions across
the program set), a bound set by the pass pool rather than the selector; variant-level
uniqueness (E7, 100\%) complements this, differentiating repeated deployments of the
same program.

For the E9 attacker study, a three-arm same-budget pool was additionally compiled:
\emph{optimized} (feature-driven), \emph{fixed-all-7} (every pass enabled, seed-only
variation, pre-registered before the runs), and \emph{Random-Pass} (uniformly random
pass application, multicompiler-style); 12 programs $\times$ 30 seeds $\times$ 2
additional arms = 720 compilation units, frozen as the attack-target pool. All three
arms reach 100\% \texttt{.text} uniqueness; Random-Pass loses 7 units on one program
(\texttt{voronoi}) to internal-compiler-error backoff (353/360 built).

\subsection{Correctness}
\label{sec:correctness}

All 12 security-subset programs pass semantic verification on all 30 seeds, and
the deployment measurements passed all 31 validation binaries. An early batch of
verification runs reported false failures on the \texttt{dc} program, traced to
concurrent verification instances sharing a scratch state file; after the
scheduler was fixed to isolate per-instance state, more than 1{,}050 subsequent
compilation--verification units completed with zero failures.

\subsection{RQ4: Deployment Cost}

\subsubsection{llama.cpp-based evaluation}
We use pp512 (prompt-processing throughput) and tg128 (token-generation throughput) on
llama.cpp running Mistral-7B-Instruct-v0.3-Q4\_K\_M-GGUF (4.07\,GB on disk, 7.25\,B
parameters), CPU-only, six threads, zero GPU layers;
three warm-ups and 20 measurements per configuration. Table~\ref{tab:llama_means}
reports the arithmetic means. The
OLLVM-full build did not produce an executable within the build-time limit
(10800.14\,s); we classify it as a build timeout, not a correctness failure, and
report no inferred runtime.

\begin{table}[!t]
\centering
\footnotesize
\caption{Arithmetic-mean llama.cpp throughput over 20 formal measurements (tokens/s).}
\label{tab:llama_means}
\begin{tabular}{lrr}
\toprule
Variant & pp512 & tg128 \\
\midrule
GCC & 31.994 & 11.276 \\
LLVM & 37.012 & 12.019 \\
optimized & 32.329 & 11.119 \\
\bottomrule
\end{tabular}
\end{table}

For pp512, optimized is 1.0\% faster than GCC but 12.7\% slower than LLVM. For tg128,
it is 1.4\% slower than GCC and 7.5\% slower than LLVM. Optimized thus remains close
to GCC in both scenarios but does not match LLVM.

\paragraph{Build time.}
Runtime is only half of the deployment cost; we also measure
full-project compilation time (Table~\ref{tab:build_time}). The complete llama.cpp
build finishes in 138.2\,s under the optimized profile against 68.5\,s under LLVM
(2.0$\times$), with balanced at 139.8\,s and GCC at 70.4\,s; the uniformly obfuscated
OLLVM-full build does not finish at all, reaching 50 of 210 compilation units (23.8\%)
before the 10{,}800\,s limit and producing no executable. On the LLVM test-suite
validation set (31 programs), the same ordering holds (optimized 10.5\,s vs.\ LLVM
6.0\,s), and the OLLVM-full build fails outright after 168.3\,s. The single-traversal
feature collection and the selection--execution decoupling are what
keep the factor near two: passes are chosen once per unit and executed from metadata
annotations, rather than each pass re-deriving its own eligibility.

\begin{table}[!t]
\centering
\footnotesize
\caption{Full-project build time. OLLVM-full on llama.cpp: timeout at
10{,}800\,s with 50/210 compilation units (23.8\%) completed, no artifact produced;
OLLVM-full on the test-suite validation set: build failure (exit 1) after 168.3\,s.}
\label{tab:build_time}
\begin{tabular}{lrr}
\toprule
Toolchain & llama.cpp (s) & test-suite, 31 programs (s) \\
\midrule
GCC & 70.4 & 5.4 \\
LLVM & 68.5 & 6.0 \\
balanced & 139.8 & 10.3 \\
optimized & 138.2 & 10.5 \\
OLLVM-full & $>$10{,}800 (timeout) & 168.3 (fail) \\
\botrule
\end{tabular}
\end{table}

\subsubsection{OpenSSL-based evaluation}
Optimized MIXCOMPILE, OLLVM-full, GCC, and LLVM are compared on AES-128-CBC,
ChaCha20-Poly1305, and SHA-512 (static build, \texttt{no-asm}, no modules; three
warm-ups and 20 one-second measurements per cell; OpenSSL~3.6.1).
Figure~\ref{fig:openssl} plots throughput across all six block sizes and
Table~\ref{tab:openssl_16384} reports the 16384-byte means.

\begin{figure}[!t]
 \centering
 \includegraphics[width=0.9\textwidth]{openssl_perf.pdf}
 \caption{(a) AES-128-CBC throughput comparison across block sizes (16\,B to 16384\,B); (b) ChaCha20-Poly1305 throughput comparison across block sizes (16\,B to 16384\,B); and (c) SHA-512 throughput comparison across block sizes (16\,B to 16384\,B).}
 \label{fig:openssl}
 \end{figure}

\begin{table*}[!t]
\centering
\footnotesize
\caption{OpenSSL arithmetic-mean throughput at 16384-byte blocks (MB/s, 20 measurements).}
\label{tab:openssl_16384}
\begin{tabular}{lrrrr}
\toprule
Algorithm & GCC & LLVM & optimized & OLLVM-full \\
\midrule
AES-128-CBC & 384.7 & 414.0 & 238.7 & 43.8 \\
ChaCha20-Poly1305 & 506.9 & 512.7 & 387.8 & 83.2 \\
SHA-512 & 685.0 & 653.0 & 241.9 & 169.4 \\
\bottomrule
\end{tabular}
\end{table*}

Across the six block sizes, optimized retains 23.3--57.9\% of LLVM throughput for
AES-128-CBC, 22.5--75.6\% for ChaCha20-Poly1305, and 20.0--37.2\% for SHA-512. At
16384 bytes, optimized reaches 57.7\%, 75.6\%, and 37.0\% of LLVM, respectively.
Optimized is slower than both mainstream compilers on every reported OpenSSL workload,
and consistently exceeds OLLVM-full at every tested algorithm and block size. The
budget gates bound long-tail expansion; the five-dimensional profile is the
instrument for trading security against throughput.

\subsubsection{Cross-compiler structural difference}
Complementing the variant-to-unprotected similarity of RQ3, this experiment measures
\emph{cross-compiler} pairwise radiff2 similarity on the 31 validation
binaries (unordered pairs among GCC, LLVM, optimized, OLLVM-full; medians over
correctness-eligible successful comparisons). Figure~\ref{fig:similarity} shows the
pairwise heatmap and Table~\ref{tab:radiff_formal} the formal medians with coverage
columns.

\begin{figure}[!t]
\centering
\includegraphics[width=0.9\textwidth]{similarity_heatmap.pdf}
\caption{Pairwise radiff2 similarity heatmap. Formal analysis uses the median of correctness-eligible successful comparisons.}
\label{fig:similarity}
\end{figure}

\begin{table*}[!t]
\centering
\footnotesize
\caption{Formal radiff2 results. The median uses correctness-eligible successful comparisons.}
\label{tab:radiff_formal}
\begin{tabular}{llrrrr}
\toprule
Variant A & Variant B & Planned & Successful & Timeouts & Eligible / Median \\
\midrule
GCC & LLVM & 31 & 31 & 0 & 31 / 0.8800 \\
GCC & optimized & 31 & 29 & 2 & 29 / 0.5620 \\
GCC & OLLVM-full & 30 & 26 & 4 & 24 / 0.1725 \\
LLVM & optimized & 31 & 29 & 2 & 29 / 0.5900 \\
LLVM & OLLVM-full & 30 & 26 & 4 & 24 / 0.1725 \\
OLLVM-full & optimized & 30 & 21 & 9 & 19 / 0.3140 \\
\bottomrule
\end{tabular}
\end{table*}

The unprotected baselines are relatively similar (GCC--LLVM median 0.880). Optimized
reduces the median to 0.562 against GCC and 0.590 against LLVM, and sits at 0.314
against OLLVM-full: selective application does not merely reproduce the structure of
uniform full obfuscation. OLLVM-full pairs have fewer eligible observations and more
timeouts, so these medians must be read together with the coverage columns.

\subsection{RQ5: Analysis Cost}

\subsubsection{Decompiler-based function recovery}
A Ghidra pipeline with ground-truth annotations (symbol-carrying LLVM reference build,
stripped target binaries, function-boundary matching) measures how much of the
original function structure an automated decompiler recovers. Recovery F1 drops from
0.903 (LLVM baseline) to 0.708 (optimized; balanced profile 0.712). This is a direct
capability metric, not a timing proxy: the automated analysis recovers substantially
less function structure from MIXCOMPILE output than from the unobfuscated build.

\subsubsection{Classical function-localization attack}
A non-LLM attacker (Ghidra plus a matching script, with full knowledge of compiler
version, pass pool, weights, and algorithm, everything except the seed, the strongest
reasonable setting) attempts to locate a target function inside the obfuscated binary.
Top-1 localization accuracy: 0.163 (optimized; balanced 0.177; unobfuscated
$\approx$1.0). This run's candidate pool differs from the E9 pool (e.g.,
\texttt{dc}: 79 candidate functions here vs.\ the frozen $\geq$10-function E9 pool);
these numbers are prior evidence that obfuscation materially increases localization
difficulty.

\subsubsection{Decompilation time}
Ghidra processes five programs under the four variants with three repetitions each
(per-program medians, then median of per-program ratios to LLVM).
Table~\ref{tab:ghidra_ratios} reports the ratios relative to LLVM (per-variant
medians: GCC 0.147\,s, LLVM 0.218\,s, OLLVM-full 0.800\,s, optimized 4.739\,s).

\begin{table}[!t]
\centering
\footnotesize
\caption{Median per-program Ghidra decompilation-time ratio relative to LLVM.}
\label{tab:ghidra_ratios}
\begin{tabular}{lr}
\toprule
Variant & Median ratio \\
\midrule
GCC & 0.677$\times$ \\
optimized & 8.109$\times$ \\
OLLVM-full & 3.626$\times$ \\
\bottomrule
\end{tabular}
\end{table}

Optimized increases the median per-program decompilation time to 8.109$\times$ LLVM,
versus 3.626$\times$ for OLLVM-full. Together with the recovery-F1 and localization
reductions above, the supported conclusion is that automated analysis of MIXCOMPILE
output is both slower and less accurate.

\subsection{Pre-Registered LLM-Era Studies: E9 and E10 Results}
\label{sec:ongoing}

Both pre-registered studies have completed under their frozen protocols. All
numbers below are final measurements; nothing is extrapolated.

\textbf{E9 (LLM-attacker arm, completed).} An LLM attacker, one frontier API model
operated in an agent harness plus a local open-weights pilot
(Qwen2.5-Coder-14B-Instruct), receives Ghidra-decompiled pseudocode of the frozen
three-arm pool (Section~\ref{sec:rq3}: A1~= fixed-all-7, A2~= optimized, plus the
unobfuscated baseline) and attempts function matching, semantic recovery, and
pass-configuration probing, across capability levels from zero-shot to
deployment-level cross-program practice (two-fold cross-validation). The
pre-registered primary endpoint was the adaptation transfer rate $\rho$: how much of
the attacker's practiced-on-other-programs improvement survives on feature-adaptive
variants versus fixed-strategy variants, a within-arm relative quantity immune to
differing obfuscation strength across arms. Four findings.
\textbf{(1)~Recognition no longer separates obfuscated from clean binaries at the
frontier tier:} the frontier attacker matches target functions at 94.4--99.2\% top-1
across arms and capability levels, against 100\% on the unobfuscated baseline
(Figure~\ref{fig:e9_layers}a,b). Protection is strongly attacker-capability
dependent: the local 14B pilot reaches only 14.2--23.3\% top-1 by arm (optimized
15.8\%), in the range of the classical localization attack (0.163).
\textbf{(2)~The surviving protection is a comprehension tax:} on the paired
semantic-recovery task ($n=39$), mean cosine to ground truth falls from 0.127
(unobfuscated) to 0.078 under fixed-all-7 ($-38.5\%$) and to 0.090 under optimized
($-28.2\%$) (Figure~\ref{fig:e9_layers}c). Obfuscation moves attacker cost from
recognition to understanding, and the optimized arm retains less of this residual
protection than the fully fixed arm. \textbf{(3)~$\rho$ is undefined:}
deployment-level practice yields no reliable gain on either arm, so the
pre-registered clause for this case applies; this attack tier does not distinguish
fixed from diversified deployment. \textbf{(4)~The pass configuration does not leak:}
configuration probing returns exact-match 0/60, with 26.7\% ``no-configuration''
answers.

\begin{figure}[!t]
\centering
\includegraphics[width=0.95\textwidth]{e9_layers.pdf}
\caption{E9: the two attack layers. Arms: A0 = unobfuscated baseline, A1 =
fixed-all-7, A2 = optimized. (a)~Frontier-attacker recognition (top-1 matching):
dots are measured cells (zero-shot and few-shot capability levels plus the
supplementary lift-to-attack cells), vertical lines span the min--max range, and the
dashed line marks the unobfuscated baseline (100\%).
(b)~Local 14B pilot recognition by arm ($n=12$ per arm); black diamonds mark the
classical localization attack's per-configuration mean.
(c)~Paired semantic-recovery cosine to ground truth ($n=39$ paired functions):
obfuscation taxes comprehension, with fixed-all-7 at $-38.5\%$ and optimized at
$-28.2\%$ relative to the unobfuscated arm.}
\label{fig:e9_layers}
\end{figure}

\textbf{E10 (LLM-guided compilation, completed).} Five arms on the same selection
task: frozen built-in selection, non-LLM search at equal budget (two sub-variants,
90/90), LLM planning without five-dimensional feedback, LLM planning with
five-dimensional feedback (invoking the interface of Section~\ref{sec:interface}),
and direct LLM source rewriting as a premise-testing arm. Two findings.
\textbf{(1)~The substrate, not the LLM, carries semantic safety:} the two LLM
planning arms jointly preserve semantics in 89/90 cases (98.9\%), their plans
executing through the compiler's constructive safeguards, while direct LLM source
rewriting preserves 43/68 (63.2\%) under the same strict verification and produces 62
non-compiling outputs besides.
\textbf{(2)~The built-in selector and LLM planners are equivalent:} the LLM planning
arms land within 1.1 percentage points of the non-LLM search arms on semantic
preservation (98.9\% vs.\ 100\%), inside the pre-registered equivalence domain
($|\Delta|\leq5$ percentage points), and the built binaries' runtime cost is governed
by the substrate rather than the planner (median ratios vs.\ the unobfuscated build:
planning-with-feedback 1.006, planning-without 1.000, rule-search 0.998,
direct-rewriting 1.002).

%=======================================================
\section{Conclusion}
%=======================================================

MIXCOMPILE turns diversified compilation into an explicit, auditable, and invocable
sampling decision: a five-dimensional cost model gives LLM-based compilation agents
a structured basis for reasoning about security--performance trade-offs, metadata
annotations decouple selection from execution behind a machine-readable interface
through which an LLM can plan, invoke, and verify compilation strategies, and
function-granularity feature extraction shapes each program's strategy by its own
structure so that recompilations remain structurally distinct. The experimental
evidence supports five statements.

First, the selection-and-cost layer that the transform literature leaves as future
work is implementable and informative: the pre-registered weights dominate 28/30
random candidates, each cost dimension has a distinct role, and random weights
silently corrupt correctness in 4/30 cases, which makes the cost model a safety
mechanism rather than decoration. Second, the mechanism's granularity is measured,
not assumed: the function--feature pairing drives per-function decisions, with
permutation moving paired Jaccard to 0.549 against a 0.708 cross-seed floor. Third,
diversity is real at both levels: 100\% variant uniqueness with program-dependent
structural distance, and between-program strategy dispersion at $3.3\times$ the
within-program seed noise, with each program's strategy stable across seeds. Fourth,
deployment cost is workload-dependent and explicitly traded: near-GCC on llama.cpp
(1--13\% below LLVM), 20.0--75.6\% of LLVM throughput on OpenSSL while consistently
outperforming full obfuscation, and a 2.0$\times$ build-time factor where uniform
full obfuscation does not finish at all. Fifth, the output measurably taxes
automated analysis: recovery F1 0.903$\to$0.708, localization top-1 0.163 vs.\
$\approx$1.0 unobfuscated, decompilation time 8.109$\times$ LLVM, and a frontier
LLM pays a 28.2\% comprehension tax where it no longer pays a recognition cost.

The evidence supports deployment when the deployed fleet is large (amortization
matters), when the workload tolerates the profiled overhead (llama.cpp-class
workloads: single-digit percent; OpenSSL-class: substantial, priced by the
five-dimensional profile), and when verification against the unobfuscated build is
part of the pipeline (the reproduction contract makes this cheap). It does not
support deployment where near-baseline throughput is non-negotiable or where
programs are known to carry undefined behavior (outside the constructive
guarantee).

MIXCOMPILE raises the per-target adaptation cost of automated and LLM-assisted
adversaries and turns each deployment into an independent draw. Security-side
validation of per-function targeting (complete vs.\ shuffled at matched overhead),
MIR-level pass pool expansion, cost-model weight learning, and deeper integration
of LLM agents with the machine-readable strategy interface are the natural
extensions.

%%%%%%%%%%%%%%%%%%%%%%%%%%%%%%%%%%%%%%%%%%%%%%%%%%%%%%%
%%% Back matter
%%%%%%%%%%%%%%%%%%%%%%%%%%%%%%%%%%%%%%%%%%%%%%%%%%%%%%%
\backmatter

\bmhead{Acknowledgements}

This work is supported by the Science and Technology Project of State Grid Corporation of China (Research on Key Technologies of the Autonomous and Controllable Power Dedicated Network Security Protection Basic Platform, $52130025002S-473-ZN$).

\section*{Declarations}

\textbf{Data and code availability.} The MIXCOMPILE framework source code (the
PassDecider selection engine and all seven obfuscation passes), the frozen
configuration files (deployment-profile and pass-vector JSON, the benchmark split
definition), and the scripts and raw measurement data underlying all reported
experiments (llama.cpp and OpenSSL throughput repetitions, radiff2 and Ghidra
decompilation results, the E1--E8 mechanism-campaign records, and the E9/E10
LLM-era study records) are available at
https://github.com/ELDARK123/MIXCOMPILE.git.

%%% Bibliography
\bibliography{sn-bibliography}

\end{document}


